\documentclass[a4paper,12pt]{article}
\usepackage{geometry}
\usepackage{hyperref}
\usepackage{tikz}
\usepackage{listings}
\usepackage{longtable}
\usepackage{booktabs}
\usepackage{graphicx}
\usepackage{xcolor}

\geometry{margin=1in}
\hypersetup{
    colorlinks=true,
    linkcolor=blue,
    filecolor=magenta,      
    urlcolor=cyan,
    pdftitle={CNSMS Project Report},
}

% Code Style
\definecolor{codegreen}{rgb}{0,0.6,0}
\definecolor{codegray}{rgb}{0.5,0.5,0.5}
\definecolor{codepurple}{rgb}{0.58,0,0.82}
\definecolor{backcolour}{rgb}{0.95,0.95,0.92}

\lstdefinestyle{pythonstyle}{
    backgroundcolor=\color{backcolour},   
    commentstyle=\color{codegreen},
    keywordstyle=\color{magenta},
    numberstyle=\tiny\color{codegray},
    stringstyle=\color{codepurple},
    basicstyle=\ttfamily\footnotesize,
    breakatwhitespace=false,         
    breaklines=true,                 
    captionpos=b,                    
    keepspaces=true,                 
    numbers=left,                    
    numbersep=5pt,                  
    showspaces=false,                
    showstringspaces=false,
    showtabs=false,                  
    tabsize=2
}
\lstset{style=pythonstyle}

\title{\textbf{Campus Navigation and Safety Management System (CNSMS)}}
\author{Muhammad Hanzala (232201020) \\ Haseeb Nawaz (232201011) \\ [Third Member] (232201015)}
\date{\today}

\begin{document}

\maketitle
\begin{abstract}
This report details the design and implementation of the Campus Navigation and Safety Management System (CNSMS). The system integrates a Python Flask backend with an Android mobile client to enhance campus security. Key features include real-time incident reporting, AI-based severity analysis, and an interactive safety heatmap. The system utilizes a hybrid authentication mechanism ensuring seamless access for both web and mobile users.
\end{abstract}

\tableofcontents
\newpage

\section{Introduction}
The CNSMS project aims to modernize campus security operations. By allowing students to report incidents via mobile devices and providing security officers with a centralized web dashboard, the system reduces response times and improves situational awareness.

\section{System Architecture}
The system follows a client-server architecture.
\begin{itemize}
    \item \textbf{Backend}: Python Flask (REST API + Web Views).
    \item \textbf{Database}: SQLite (Storage).
    \item \textbf{Clients}: Android App (Student View) and Web Dashboard (Admin View).
\end{itemize}

\begin{figure}[h]
    \centering
    \begin{tikzpicture}
        \node[draw, rectangle, minimum width=3cm, minimum height=1.5cm] (Android) at (0,0) {Android App};
        \node[draw, rectangle, minimum width=3cm, minimum height=1.5cm] (Web) at (0,3) {Web Dashboard};
        \node[draw, rectangle, minimum width=4cm, minimum height=3cm] (Server) at (6,1.5) {Flask Server};
        \node[draw, cylinder, shape border rotate=90, aspect=0.25, minimum height=2cm, minimum width=1.5cm] (DB) at (10,1.5) {SQLite DB};
        
        \draw[->, thick] (Android) -- node[above] {REST API (JWT)} (Server);
        \draw[->, thick] (Web) -- node[above] {HTTP (Session)} (Server);
        \draw[<->, thick] (Server) -- (DB);
    \end{tikzpicture}
    \caption{High-Level System Architecture}
\end{figure}

\section{Backend Implementation}
The backend is structured using Flask Blueprints for modularity.

\subsection{Authentication}
A hybrid approach handles differing client needs:
\begin{itemize}
    \item \textbf{Web}: Uses \texttt{Flask-Login} for session management.
    \item \textbf{API}: Uses \texttt{PyJWT} for stateless token authentication.
\end{itemize}

\subsection{AI Integration}
The system includes an AI service module designed to analyze uploaded incident images.
\begin{lstlisting}[language=Python, caption=AI Analysis Logic Stub]
def analyze_image(image_bytes):
    # Mock analysis logic
    return {
        "severity": "HIGH",
        "confidence": 0.88,
        "labels": ["fire", "hazard"]
    }
\end{lstlisting}

\section{Database Schema}
The database manages users, incidents, and audit logs.

\begin{longtable}{|p{3cm}|p{4cm}|p{6cm}|}
\hline
\textbf{Table} & \textbf{Key Columns} & \textbf{Description} \\
\hline
\endhead
User & id, email, role & Stores credentials and access levels. \\
\hline
Incident & id, user\_id, status & Stores incident reports and AI results. \\
\hline
\end{longtable}

\section{Security & Usability}
\subsection{Security Measures}
\begin{itemize}
    \item Password hashing with Werkzeug (PBKDF2).
    \item JWT tokens for API protection.
    \item Input validation to prevent injection.
\end{itemize}

\subsection{HCI Features}
\begin{itemize}
    \item Responsive Dashboard (Bootstrap 5).
    \item Color-coded status indicators (Red=High Risk).
    \item Toast notifications for user feedback.
\end{itemize}

\section{Conclusion}
The CNSMS prototype successfully demonstrates the core flow of incident reporting and management. Future work includes integrating a live AI model and deploying to a cloud provider.

\end{document}
